% ═══════════════════════════════════════════════════════════════
%  ch03.tex — فصل ۳
%  مبانی حقوقی مداخله‌ی خارجی در حقوق بین‌الملل
%  نسخه‌ی اصلاح‌شده: رفع مشکلات RTL، رنگ و جهت
% ═══════════════════════════════════════════════════════════════

\chapter{مبانی حقوقی مداخله‌ی خارجی در حقوق بین‌الملل}
\label{ch:legal}

\begin{tcolorbox}[quotebox, title={}]
	«تمام ملت‌ها حق تعیین سرنوشت دارند. بر اساس این حق، آن‌ها
	آزادانه وضعیت سیاسی خود را تعیین و آزادانه توسعه‌ی اقتصادی،
	اجتماعی و فرهنگی خود را دنبال می‌کنند.»
	
	\hfill ---ماده‌ی ۱ مشترک میثاقین بین‌المللی
	۱۹۶۶\src{\lr{ICCPR \& ICESCR, Art.\,1}}
\end{tcolorbox}


% ────────────────────────────────────────────────────────────────
\section{درآمد}
\label{sec:ch03-intro}

فصل‌های پیشین چارچوب مفهومی و فلسفی مداخله‌ی خارجی را ترسیم
کردند. اما مداخله‌ی خارجی صرفاً مسئله‌ای فلسفی نیست؛ ابعاد
حقوقی دقیقی دارد که در حقوق بین‌الملل معاصر تدوین شده‌اند.
این فصل به بررسی سه محور حقوقی اصلی می‌پردازد:
\begin{enumerate}[label=\textbf{\arabic*.}, nosep, rightmargin=1em]
	\item اصل \concept{حاکمیت ملی} و \concept{عدم مداخله}؛
	\item حق \concept{تعیین سرنوشت} و حق \concept{جدایی جبرانی}؛
	\item دکترین \concept{مسئولیت حمایت} (\lr{R2P}).
\end{enumerate}


% ────────────────────────────────────────────────────────────────
\section{اصل حاکمیت ملی و عدم مداخله}
\label{sec:ch03-sovereignty}

\subsection{از وستفالی تا منشور ملل متحد}
\label{subsec:ch03-westphalia}

\begin{tcolorbox}[defbox,
	title={مفهوم: حاکمیت وستفالیایی}]
	\concept{حاکمیت وستفالیایی}%
	\footnote{حاکمیت وستفالیایی؛ به انگلیسی:
		\lr{Westphalian Sovereignty}}%
	اصلی است برخاسته از معاهدات صلح وستفالی (۱۶۴۸) که بر
	دو رکن استوار است:
	\begin{enumerate}[label=\textbf{\alph*.}, nosep, rightmargin=1em]
		\item \concept{سرزمینی بودن}: هر دولت بر سرزمین مشخصی
		اقتدار انحصاری دارد؛
		\item \concept{عدم مداخله}: هیچ دولتی حق مداخله در امور
		داخلی دولت دیگر را ندارد.
	\end{enumerate}
	این اصل در ماده‌ی ۲(۴) و ماده‌ی ۲(۷) منشور ملل متحد (۱۹۴۵)
	تدوین رسمی یافته
	است\src{\lr{UN Charter, Art.\,2(4), 2(7)}}.
\end{tcolorbox}

\thinker{استفن کراسنر}%
\footnote{%
	استفن کراسنر (\lr{Stephen D.\,Krasner, 1942–})؛%
	استاد علوم سیاسی دانشگاه استنفورد.%
}%
در اثر مهم خود
(\lr{\textit{Sovereignty: Organized Hypocrisy}}، ۱۹۹۹)
استدلال می‌کند که حاکمیت وستفالیایی همواره بیش از آنکه واقعیت
باشد، \concept{هنجار آرمانی} بوده
است\src{\lr{Krasner, 1999, pp.\,3--42}}.
دولت‌ها در عمل بارها و بارها در امور یکدیگر مداخله کرده‌اند؛
حتی در دوره‌ای که ادعا می‌شد اصل عدم مداخله مطلق است.

\subsection{استثنائات حقوقی اصل عدم مداخله}
\label{subsec:ch03-exceptions}

منشور ملل متحد دو استثناء رسمی بر اصل عدم مداخله پیش‌بینی
کرده است:

\begin{enumerate}[label=\textbf{\arabic*.}, nosep, rightmargin=1em]
	\item \concept{دفاع مشروع فردی یا دسته‌جمعی}
	(ماده‌ی ۵۱ منشور)\src{\lr{UN Charter, Art.\,51}}: هرگاه
	حمله‌ی مسلحانه‌ای صورت گیرد، دولت قربانی حق دفاع دارد
	و می‌تواند از دیگران کمک بخواهد؛
	\item \concept{اقدام شورای امنیت تحت فصل هفتم}
	(مواد ۳۹ تا ۴۲ منشور)\src{\lr{UN Charter, Ch.\,VII}}: شورای
	امنیت می‌تواند تهدید علیه صلح، نقض صلح یا عمل تجاوز را
	احراز و اقدامات اجباری، از جمله نظامی، را تجویز کند.
\end{enumerate}

\begin{tcolorbox}[legalbox,
	title={خلأ حقوقی: مداخله بدون مجوز شورای امنیت}]
	مسئله‌ی اصلی هنگامی ظاهر می‌شود که یک یا چند عضو دائم شورای
	امنیت (با حق وتو) مانع صدور مجوز مداخله شوند، حتی در مواجهه
	با نسل‌کشی. نمونه‌ی بارز: قتل‌عام رواندا (۱۹۹۴) که شورای امنیت
	به دلایل سیاسی از مداخله خودداری
	کرد\src{\lr{Dallaire, 2003, pp.\,264--299}}.
	این خلأ حقوقی موجب طرح دکترین‌های جایگزین شده که مهم‌ترین آن‌ها
	\concept{مسئولیت حمایت} است.
\end{tcolorbox}


% ────────────────────────────────────────────────────────────────
\section{حق تعیین سرنوشت}
\label{sec:ch03-self-determination}

\subsection{تحول تاریخی مفهوم}
\label{subsec:ch03-sd-history}

\begin{figure}[htbp]
	\centering
	\begin{tikzpicture}[
		>=Stealth,
		every node/.style={font=\small},
		evt/.style={
			draw=PrimaryMid, fill=BgCool, rounded corners=3pt,
			text width=2.8cm, align=center, minimum height=1.4cm,
			font=\small
		}
		]
		% خط زمانی
		\draw[->, very thick, PrimaryDeep] (0,0) -- (15.5,0);
		\node[below, font=\footnotesize] at (15,-0.2) {\rl{زمان}};
		
		% رویدادها — صریح (بدون foreach)
		\node[evt] (e1) at (1.5,2)
		{\textbf{\rl{۱۷۷۶/۱۷۸۹}}\\[2pt]%
			\rl{انقلاب آمریکا}\\[1pt]%
			\rl{و فرانسه}};
		
		\node[evt] (e2) at (4.5,2)
		{\textbf{\rl{۱۹۱۸}}\\[2pt]%
			\rl{اصول چهارده‌گانه}\\[1pt]%
			\rl{ویلسون}};
		
		\node[evt] (e3) at (7.5,2)
		{\textbf{\rl{۱۹۴۵}}\\[2pt]%
			\rl{منشور ملل متحد}\\[1pt]%
			\rl{مواد ۱ و ۵۵}};
		
		\node[evt] (e4) at (10.5,2)
		{\textbf{\rl{۱۹۶۰--۷۰}}\\[2pt]%
			\rl{استعمارزدایی}\\[1pt]%
			\rl{قطعنامه ۱۵۱۴}};
		
		\node[evt] (e5) at (14,2)
		{\textbf{\rl{۱۹۹۰--اکنون}}\\[2pt]%
			\rl{جدایی جبرانی}\\[1pt]%
			\rl{کوزوو، سودان ج.}};
		
		% خطوط اتصال — صریح
		\draw[thick, AccentGold] (e1.south) -- (1.5,0);
		\draw[thick, AccentGold] (e2.south) -- (4.5,0);
		\draw[thick, AccentGold] (e3.south) -- (7.5,0);
		\draw[thick, AccentGold] (e4.south) -- (10.5,0);
		\draw[thick, AccentGold] (e5.south) -- (14,0);
		
		% عنوان
		\node[above, font=\normalsize\bfseries, text=PrimaryDeep]
		at (7.5,4) {\rl{تحول تاریخی حق تعیین سرنوشت}};
	\end{tikzpicture}
	\caption{خط زمانی تحول مفهوم حق تعیین سرنوشت}
	\label{fig:ch03-sd-timeline}
\end{figure}

حق تعیین سرنوشت از یک ایده‌ی فلسفی (انقلاب فرانسه) به یک
اصل حقوقی الزام‌آور (میثاقین ۱۹۶۶) تحول یافته
است\src{\lr{Cassese, 1995, pp.\,11--33}}.
\thinker{آنتونیو کاسسه}%
\footnote{%
	آنتونیو کاسسه (\lr{Antonio Cassese, 1937--2011})؛%
	حقوقدان بین‌المللی ایتالیایی و نخستین رئیس دادگاه کیفری
	بین‌المللی برای یوگسلاوی سابق.%
}%
تمایز مهمی میان دو بعد این حق قائل شده است:

\begin{tcolorbox}[comparebox,
	title={دو بعد حق تعیین سرنوشت}]
	\begin{tabular}{>{\bfseries\color{PrimaryMid}}p{3cm}
			p{5cm} p{5cm}}
		\toprule
		\textbf{بعد} &
		\textbf{تعیین سرنوشت خارجی} &
		\textbf{تعیین سرنوشت داخلی} \\
		\midrule
		تعریف &
		حق تشکیل دولت مستقل یا الحاق به دولت دیگر &
		حق مشارکت در حکمرانی و خودمختاری درون دولت موجود \\
		\midrule
		مصداق کلاسیک &
		استعمارزدایی (قطعنامه‌ی ۱۵۱۴) &
		خودمختاری قومی، فدرالیسم \\
		\midrule
		وضعیت حقوقی &
		پذیرفته‌شده در بافت استعمار؛ مناقشه‌آمیز خارج از آن &
		به‌طور گسترده پذیرفته‌شده \\
		\midrule
		رابطه با مداخله &
		ممکن است مداخله‌ی خارجی را مشروع سازد &
		عمدتاً امر داخلی است \\
		\bottomrule
	\end{tabular}
\end{tcolorbox}

\subsection{جدایی جبرانی: بسط حقوقی}
\label{subsec:ch03-remedial}

\begin{tcolorbox}[legalbox,
	title={نظریه‌ی جدایی جبرانی در حقوق بین‌الملل}]
	نظریه‌ی \concept{جدایی جبرانی}%
	\footnote{جدایی جبرانی؛ به انگلیسی: \lr{Remedial Secession}}%
	بر اساس رأی مشورتی دیوان بین‌المللی دادگستری درباره‌ی کوزوو
	(۲۰۱۰) و تفسیر بند حفاظتی قطعنامه‌ی ۲۶۲۵ مجمع عمومی (۱۹۷۰)
	استدلال می‌کند که:
	
	اگر دولتی \textbf{حق تعیین سرنوشت داخلی} گروهی از شهروندانش
	را به‌طور سیستماتیک و مستمر نقض کند و هیچ ابزار اصلاحی دیگری
	مؤثر نباشد، آن گروه حق تعیین سرنوشت \textbf{خارجی} (جدایی)
	را به‌صورت \concept{جبرانی} به دست
	می‌آورد\src{\lr{ICJ, Advisory Opinion on Kosovo, 2010;
			Buchanan, 2004, pp.\,331--400; Tancredi, 2006, pp.\,171--200}}.
\end{tcolorbox}

با این حال، هنوز اجماع حقوقی کاملی درباره‌ی جدایی جبرانی وجود
ندارد. بسیاری از دولت‌ها، به‌ویژه دولت‌هایی که خود با جنبش‌های
جدایی‌طلبانه مواجه‌اند (چین، روسیه، هند، ترکیه، ایران، اسپانیا)
مخالف شناسایی این حق هستند.


% ────────────────────────────────────────────────────────────────
\section{دکترین مسئولیت حمایت (\lr{R2P})}
\label{sec:ch03-r2p}

\subsection{پیشینه و شکل‌گیری}
\label{subsec:ch03-r2p-history}

شکست جامعه‌ی بین‌المللی در جلوگیری از نسل‌کشی رواندا (۱۹۹۴) و
قتل‌عام سربرنیتسا (۱۹۹۵) بحران عمیقی در حقوق بین‌الملل ایجاد
کرد. از یک سو، اصل حاکمیت ملی مانع مداخله شده بود؛ از سوی دیگر،
نتیجه‌ی عدم مداخله کشتار صدها هزار غیرنظامی بود.

\thinker{کوفی عنان}%
\footnote{%
	کوفی عنان (\lr{Kofi Annan, 1938--2018})؛%
	هفتمین دبیرکل سازمان ملل متحد.%
}%
در سال ۱۹۹۹ پرسش بنیادینی مطرح کرد: «اگر مداخله‌ی بشردوستانه
واقعاً حمله‌ی غیرقابل قبولی به حاکمیت است، پس ما به مردمی که
در رواندا و سربرنیتسا کشته شدند چه
می‌گوییم؟»\src{\lr{Annan, 1999}}

در پاسخ به این چالش، دولت کانادا در سال ۲۰۰۰
\concept{کمیسیون بین‌المللی مداخله و حاکمیت ملی}%
\footnote{کمیسیون بین‌المللی مداخله و حاکمیت ملی؛ به انگلیسی:
	\lr{International Commission on Intervention and State
		Sovereignty (ICISS)}.}%
را تشکیل داد. گزارش این کمیسیون، با عنوان
«مسئولیت حمایت» (۲۰۰۱)، نقطه‌ی
عطفی در حقوق بین‌الملل بود\src{\lr{ICISS, 2001}}.

\subsection{سه رکن مسئولیت حمایت}
\label{subsec:ch03-r2p-pillars}

\begin{tcolorbox}[legalbox,
	title={سه رکن دکترین مسئولیت حمایت}]
	
	\needspace{4\baselineskip}
	\subsubsection{رکن اول: مسئولیت دولت}
	هر دولت مسئولیت اولیه‌ی حمایت از جمعیت خود را در برابر
	\concept{نسل‌کشی}، \concept{جنایات جنگی}، \concept{پاکسازی
		قومی} و \concept{جنایات علیه بشریت} بر عهده
	دارد\src{\lr{UN World Summit Outcome, 2005, para.\,138}}.
	
	\needspace{4\baselineskip}
	\subsubsection{رکن دوم: مسئولیت جامعه‌ی بین‌المللی در کمک}
	جامعه‌ی بین‌المللی باید به دولت‌ها در انجام این مسئولیت کمک
	کند؛ از طریق ظرفیت‌سازی، هشدار زودهنگام و دیپلماسی
	پیشگیرانه\src{\lr{ibid., para.\,138}}.
	
	\needspace{4\baselineskip}
	\subsubsection{رکن سوم: مسئولیت واکنش}
	اگر دولتی آشکارا از حمایت جمعیت خود ناتوان یا ناخواه باشد،
	جامعه‌ی بین‌المللی مسئولیت اقدام جمعی، از جمله اقدام نظامی
	به‌عنوان آخرین چاره و از طریق شورای امنیت، را بر عهده
	دارد\src{\lr{ibid., para.\,139}}.
\end{tcolorbox}

\begin{figure}[htbp]
	\centering
	\begin{tikzpicture}[
		>=Stealth,
		every node/.style={font=\small},
		pillar/.style={
			draw=PrimaryMid, fill=#1, rounded corners=4pt,
			minimum width=3.8cm, minimum height=1.8cm,
			text width=3.4cm, align=center, font=\small
		}
		]
		% سه ستون
		\node[pillar=BgTeal!50] (p1) at (0,0)
		{\textbf{\rl{رکن اول}}\\[3pt]%
			\rl{مسئولیت دولت}\\[1pt]%
			{\scriptsize\rl{حمایت از جمعیت خود}}};
		
		\node[pillar=BgWarm!60] (p2) at (5.5,0)
		{\textbf{\rl{رکن دوم}}\\[3pt]%
			\rl{کمک بین‌المللی}\\[1pt]%
			{\scriptsize\rl{ظرفیت‌سازی و پیشگیری}}};
		
		\node[pillar=red!10] (p3) at (11,0)
		{\textbf{\rl{رکن سوم}}\\[3pt]%
			\rl{واکنش جمعی}\\[1pt]%
			{\scriptsize\rl{مداخله: آخرین چاره}}};
		
		% فلش‌ها
		\draw[->, very thick, AccentGold] (p1.east) -- (p2.west)
		node[midway, above, font=\scriptsize, text=AccentAmber]
		{\rl{ناتوانی دولت}};
		\draw[->, very thick, AccentRed] (p2.east) -- (p3.west)
		node[midway, above, font=\scriptsize, text=AccentRed]
		{\rl{شکست پیشگیری}};
		
		% پایه
		\draw[very thick, PrimaryDeep]
		(-2.5,-1.5) -- (13.5,-1.5);
		\node[below, font=\normalsize\bfseries, text=PrimaryDeep]
		at (5.5,-1.8) {\rl{دکترین مسئولیت حمایت (\lr{R2P})}};
	\end{tikzpicture}
	\caption{سه رکن دکترین مسئولیت حمایت}
	\label{fig:ch03-r2p-pillars}
\end{figure}

\subsection{اعمال \lr{R2P} در عمل}
\label{subsec:ch03-r2p-practice}

\begin{table}[htbp]
	\centering
	\caption{موارد اعمال یا عدم اعمال مسئولیت حمایت}
	\label{tab:ch03-r2p-cases}
	\begin{adjustbox}{max width=\textwidth}
		\begin{tabular}{>{\bfseries\color{PrimaryMid}\small}p{2cm}
				>{\small}p{2.5cm} >{\small}p{2.2cm}
				>{\small}p{2.5cm} >{\small}p{3cm}}
			\toprule
			\textbf{مورد} & \textbf{بحران} & \textbf{تصمیم ش.ا.} &
			\textbf{اقدام} & \textbf{نتیجه} \\
			\midrule
			کنیا ۲۰۰۸ &
			خشونت پس از انتخابات &
			دیپلماسی &
			میانجی‌گری عنان &
			موفق: توقف خشونت \\
			\midrule
			لیبی ۲۰۱۱ &
			سرکوب معترضان &
			قطعنامه‌ی ۱۹۷۳ &
			حمله‌ی هوایی ناتو &
			سقوط قذافی؛ فروپاشی دولت \\
			\midrule
			ساحل عاج ۲۰۱۱ &
			بحران پس از انتخابات &
			قطعنامه‌ی ۱۹۷۵ &
			مداخله‌ی فرانسه و س.م. &
			نسبتاً موفق \\
			\midrule
			سوریه ۲۰۱۱-- &
			جنگ داخلی و جنایات &
			وتوی روسیه و چین &
			عدم مداخله‌ی رسمی &
			فاجعه‌ی انسانی \\
			\midrule
			میانمار ۲۰۱۷-- &
			پاکسازی روهینگیا &
			بیانیه‌ی ضعیف &
			عدم اقدام مؤثر &
			تداوم بحران \\
			\bottomrule
		\end{tabular}
	\end{adjustbox}
\end{table}

\subsection{نقد مسئولیت حمایت}
\label{subsec:ch03-r2p-critique}

\begin{tcolorbox}[enemybox,
	title={نقدهای اصلی بر دکترین \lr{R2P}}]
	\begin{enumerate}[label=\textbf{\arabic*.}, nosep, rightmargin=1em]
		\item \enemy{گزینشی بودن}: چرا لیبی بله و سوریه نه؟ منتقدان
		استدلال می‌کنند که \lr{R2P} تنها علیه دولت‌های ضعیف و بدون
		حامی قدرتمند اعمال
		می‌شود\src{\lr{Hehir, 2012, pp.\,68--95}}؛
		\item \enemy{سوءاستفاده}: تجربه‌ی لیبی نشان داد که مجوز
		«حمایت از غیرنظامیان» به عملیات تغییر رژیم تبدیل شد.
		روسیه و چین این تجربه را دلیل مخالفت با هرگونه قطعنامه‌ی
		مشابه درباره‌ی سوریه قرار
		دادند\src{\lr{Bellamy, 2015, pp.\,159--178}}؛
		\item \enemy{نو-استعمارگرایی}: برخی دولت‌های جنوب جهانی
		\lr{R2P} را پوشش حقوقی جدیدی برای مداخله‌ی استعماری تلقی
		می‌کنند\src{\lr{Mamdani, 2009, pp.\,49--72}}؛
		\item \enemy{ناکارآمدی ساختاری}: وابستگی به شورای امنیت و
		حق وتو، دکترین را در مهم‌ترین موارد فلج می‌کند؛
		\item \enemy{غفلت از پیشگیری}: تمرکز بر رکن سوم (واکنش نظامی)
		به حاشیه‌رانی رکن‌های اول و دوم انجامیده
		است\src{\lr{Evans, 2008, pp.\,56--72}}.
	\end{enumerate}
\end{tcolorbox}


% ────────────────────────────────────────────────────────────────
\section{مداخله به دعوت: مبانی حقوقی}
\label{sec:ch03-invitation}

\subsection{تعریف و شرایط حقوقی}
\label{subsec:ch03-inv-def}

\begin{tcolorbox}[legalbox,
	title={مداخله به دعوت در حقوق بین‌الملل}]
	\concept{مداخله به دعوت}%
	\footnote{مداخله به دعوت؛ به انگلیسی:
		\lr{Intervention by Invitation}}%
	وضعیتی است که در آن دولت مشروع یک کشور، رسماً از یک یا چند
	دولت خارجی درخواست مداخله‌ی نظامی در خاک خود می‌کند. این شکل
	از مداخله در حقوق بین‌الملل عرفی عمدتاً مجاز شناخته
	شده\src{\lr{Doswald-Beck, 1985, pp.\,189--249}}
	ولی شرایطی دارد:
	\begin{enumerate}[label=\textbf{\alph*.}, nosep, rightmargin=1em]
		\item دعوت باید از \concept{دولت مشروع} باشد (نه از شورشیان)؛
		\item دعوت باید \concept{آزادانه} و بدون اجبار باشد؛
		\item مداخله باید \concept{متناسب} با درخواست باشد؛
		\item مداخله نباید ناقض \concept{حق تعیین سرنوشت} مردم باشد.
	\end{enumerate}
\end{tcolorbox}

\subsection{ابهامات و چالش‌ها}
\label{subsec:ch03-inv-challenges}

\thinker{لوئیز داسپرمون}%
\footnote{%
	ژان داسپرمون (\lr{Jean d'Aspremont})؛%
	استاد حقوق بین‌الملل دانشگاه منچستر.%
}%
چهار ابهام اساسی در مفهوم «مداخله به دعوت» شناسایی کرده
است\src{\lr{d'Aspremont, 2015, pp.\,24--68}}:

\begin{enumerate}[label=\textbf{\arabic*.}, nosep, rightmargin=1em]
	\item \concept{بحران مشروعیت دعوت‌کننده}: در جنگ داخلی، کدام
	طرف «دولت مشروع» است؟
	\item \concept{رضایت واقعی یا صوری}: آیا دعوت تحت فشار انجام
	شده؟
	\item \concept{گستره‌ی مداخله}: دعوت برای یک هدف مشخص ممکن
	است به مداخله‌ی گسترده تبدیل شود؛
	\item \concept{حق تعیین سرنوشت}: اگر بخش بزرگی از مردم با
	دولت دعوت‌کننده مخالف باشند، آیا دعوت نقض حق تعیین سرنوشت
	نیست؟
\end{enumerate}

\begin{tcolorbox}[casebox,
	title={مطالعه‌ی موردی: مداخلات «دعوت‌شده» در جنگ سرد}]
	\begin{table}[H]
		\centering
		\begin{adjustbox}{max width=\linewidth}
			\begin{tabular}{>{\bfseries\color{PrimaryMid}\small}p{2.5cm}
					>{\small}p{2.8cm} >{\small}p{2.5cm} >{\small}p{3.5cm}}
				\toprule
				\textbf{مورد} & \textbf{دعوت‌کننده} & \textbf{مداخله‌گر} &
				\textbf{واقعیت} \\
				\midrule
				مجارستان ۱۹۵۶ &
				دولت دست‌نشانده &
				شوروی &
				سرکوب انقلاب مردمی \\
				\midrule
				چکسلواکی ۱۹۶۸ &
				«دعوت» جعلی &
				شوروی و پیمان ورشو &
				پایان بهار پراگ \\
				\midrule
				افغانستان ۱۹۷۹ &
				دولت کودتایی &
				شوروی &
				اشغال ده‌ساله \\
				\midrule
				گرانادا ۱۹۸۳ &
				فرماندار کل (نمادین) &
				آمریکا &
				تغییر رژیم \\
				\midrule
				پاناما ۱۹۸۹ &
				رئیس‌جمهور منتخب (تبعیدی) &
				آمریکا &
				دستگیری نوریگا \\
				\bottomrule
			\end{tabular}
		\end{adjustbox}
	\end{table}
\end{tcolorbox}

% ═══════════════════════════════════════════════════════════════
%  ch03.tex — ادامه (از بخش دعوت غیردولتی)
% ═══════════════════════════════════════════════════════════════

% ────────────────────────────────────────────────────────────────
\section{دعوت از منظر غیردولتی: اقلیت‌ها و جنبش‌های آزادی‌بخش}
\label{sec:ch03-non-state}

\subsection{جایگاه حقوقی جنبش‌های آزادی‌بخش}
\label{subsec:ch03-liberation}

در دهه‌های ۱۹۶۰ و ۱۹۷۰، مجمع عمومی سازمان ملل متحد به
\concept{جنبش‌های آزادی‌بخش ملی}%
\footnote{جنبش‌های آزادی‌بخش ملی؛ به انگلیسی:
	\lr{National Liberation Movements (NLMs)}}%
در بافت استعمارزدایی مشروعیت بخشید. قطعنامه‌ی ۳۱۰۳ (۱۹۷۳)
حق مبارزه‌ی مسلحانه‌ی مردمان تحت استعمار و اشغال را به رسمیت
شناخت\src{\lr{UNGA Res.\,3103 (XXVIII), 1973}}.

اما آیا این حق به فراتر از بافت استعمارزدایی قابل تعمیم است؟
آیا یک اقلیت قومی تحت ستم، مانند کردهای عراق، اویغورهای
چین یا ایزدی‌های تحت اشغال داعش، نیز حق دعوت از نیروی
خارجی را دارند؟

\begin{tcolorbox}[comparebox,
	title={دعوت دولتی در برابر دعوت غیردولتی}]
	\begin{tabular}{>{\bfseries\color{PrimaryMid}}p{2.8cm}
			p{5cm} p{5cm}}
		\toprule
		\textbf{معیار} &
		\textbf{دعوت دولتی} &
		\textbf{دعوت غیردولتی (اقلیت/جنبش)} \\
		\midrule
		مبنای حقوقی &
		حاکمیت ملی و حقوق بین‌الملل عرفی &
		حق تعیین سرنوشت و حقوق بشر \\
		\midrule
		مشروعیت عمومی &
		نسبتاً پذیرفته‌شده &
		بسیار مناقشه‌آمیز \\
		\midrule
		ریسک سوءاستفاده &
		متوسط &
		بالا \\
		\midrule
		نمونه‌ی تاریخی &
		سوریه-روسیه (۲۰۱۵) &
		کردهای عراق-آمریکا (۱۹۹۱/۲۰۱۴) \\
		\midrule
		چالش اصلی &
		مشروعیت دولت دعوت‌کننده &
		نمایندگی جنبش و عواقب بلندمدت \\
		\bottomrule
	\end{tabular}
\end{tcolorbox}


% ────────────────────────────────────────────────────────────────
\section{حقوق بین‌الملل بشردوستانه و حمایت از اقلیت‌ها}
\label{sec:ch03-humanitarian}

\subsection{چارچوب کنوانسیونی}
\label{subsec:ch03-conventions}

مجموعه‌ای از اسناد بین‌المللی حمایت از گروه‌های آسیب‌پذیر و
اقلیت‌ها را تضمین می‌کنند:

\begin{enumerate}[label=\textbf{\arabic*.}, nosep, rightmargin=1em]
	\item \concept{کنوانسیون منع و مجازات نسل‌کشی} (۱۹۴۸):
	دولت‌ها متعهد به پیشگیری و مجازات نسل‌کشی‌اند، نه فقط
	در قلمرو خود بلکه در هر
	جا\src{\lr{Genocide Convention, 1948, Art.\,I}}؛
	\item \concept{کنوانسیون‌های چهارگانه‌ی ژنو} (۱۹۴۹) و
	پروتکل‌های الحاقی: حمایت از غیرنظامیان در جنگ؛
	\item \concept{ماده‌ی ۲۷ میثاق بین‌المللی حقوق مدنی و سیاسی}:
	حقوق اقلیت‌های قومی، زبانی و
	مذهبی\src{\lr{ICCPR, Art.\,27}}؛
	\item \concept{اعلامیه‌ی حقوق اشخاص متعلق به اقلیت‌ها}
	(۱۹۹۲)\src{\lr{UNGA Res.\,47/135, 1992}}؛
	\item \concept{اساسنامه‌ی رم} (۱۹۹۸): تعریف جنایات علیه
	بشریت و جنایات جنگی و ایجاد دیوان کیفری بین‌المللی.
\end{enumerate}

\subsection{شکاف میان هنجار و واقعیت}
\label{subsec:ch03-gap}

\begin{tcolorbox}[enemybox,
	title={شکاف هنجار-واقعیت در حمایت از اقلیت‌ها}]
	علیرغم وجود این چارچوب حقوقی گسترده، واقعیت آن است که:
	\begin{itemize}[nosep, rightmargin=1em]
		\item اویغورها در اردوگاه‌های بازآموزی چین
		هستند\src{\lr{Zenz, 2019, pp.\,102--128}}؛
		\item روهینگیاها از میانمار پاکسازی قومی
		شدند\src{\lr{UN FFM on Myanmar, 2018}}؛
		\item ایزدی‌ها نسل‌کشی داعش را تجربه
		کردند\src{\lr{UN HRC, 2016, A/HRC/32/CRP.2}}؛
		\item چچنی‌ها دو جنگ ویرانگر را از سر
		گذراندند\src{\lr{Politkovskaya, 2003, pp.\,1--42}}؛
		\item گرجی‌ها و آذری‌ها در قرن نوزدهم مورد ستم سیستماتیک
		قاجار و عثمانی قرار
		گرفتند\src{\lr{Suny, 1994, pp.\,28--63}}.
	\end{itemize}
	در اغلب این موارد، جامعه‌ی بین‌المللی، علیرغم داشتن
	ابزارهای حقوقی، به دلایل سیاسی از اقدام مؤثر بازمانده است.
\end{tcolorbox}


% ────────────────────────────────────────────────────────────────
\section{استمداد از خارج برای رفع استبداد: مرور حقوقی}
\label{sec:ch03-appeal}

\subsection{آیا مردم حق دارند از بیگانه کمک بخواهند؟}
\label{subsec:ch03-appeal-right}

این پرسش یکی از خلأهای بزرگ حقوق بین‌الملل است. حقوق بین‌الملل
موضع روشنی درباره‌ی «حق استمداد»
(\lr{Right to Appeal for External Assistance}) ندارد.
آنچه وجود دارد:

\begin{enumerate}[label=\textbf{\arabic*.}, nosep, rightmargin=1em]
	\item حق \concept{شکایت فردی} نزد نهادهای بین‌المللی حقوق بشری
	(کمیته‌ی حقوق بشر، دیوان اروپایی حقوق بشر)؛
	\item حق \concept{پناهندگی} (کنوانسیون ۱۹۵۱)؛
	\item مشروعیت عرفی \concept{درخواست کمک بشردوستانه}
	(غیرنظامی)؛
	\item مشروعیت مناقشه‌آمیز \concept{درخواست مداخله‌ی نظامی}.
\end{enumerate}

\thinker{آن‌ماری اسلاتر}%
\footnote{%
	آن‌ماری اسلاتر (\lr{Anne-Marie Slaughter, 1958–})؛%
	استاد حقوق بین‌الملل دانشگاه پرینستون.%
}%
استدلال می‌کند که در عصر جهانی‌شدن، مفهوم حاکمیت باید از
«حاکمیت به‌مثابه‌ی کنترل» به «حاکمیت به‌مثابه‌ی مسئولیت»
تحول یابد\src{\lr{Slaughter, 2004, pp.\,261--283}}.
در این چارچوب، دولتی که مسئولیت خود را در قبال شهروندانش ایفا
نکند، نمی‌تواند به حاکمیت خود برای جلوگیری از کمک خارجی
استناد کند.

\subsection{خط قرمز: مداخله‌ی بشردوستانه در برابر تغییر رژیم}
\label{subsec:ch03-regime-change}

\begin{tcolorbox}[legalbox,
	title={تمایز حقوقی: حمایت از غیرنظامیان و تغییر رژیم}]
	حقوق بین‌الملل معاصر تمایز مهمی میان دو هدف قائل است:
	\begin{enumerate}[label=\textbf{\alph*.}, nosep, rightmargin=1em]
		\item \concept{حمایت از غیرنظامیان}
		(\lr{Protection of Civilians}): مجاز بر اساس فصل هفتم
		منشور و \lr{R2P}؛
		\item \concept{تغییر رژیم}
		(\lr{Regime Change}): غیرمجاز و مغایر با اصل حاکمیت
		ملی\src{\lr{Reisman, 2000, pp.\,79--94}}.
	\end{enumerate}
	
	مشکل اصلی آن است که در عمل، این دو اغلب درهم‌آمیخته‌اند.
	قطعنامه‌ی ۱۹۷۳ درباره‌ی لیبی (۲۰۱۱) «حمایت از غیرنظامیان» را
	مجاز کرد، اما عملیات ناتو عملاً به تغییر رژیم و قتل قذافی
	انجامید. این تجربه اعتماد به \lr{R2P} را به‌شدت تضعیف
	کرد\src{\lr{Kuperman, 2013, pp.\,105--148}}.
\end{tcolorbox}


% ────────────────────────────────────────────────────────────────
\section{جدایی و تشکیل کشور جدید: نتایج متفاوت}
\label{sec:ch03-new-states}

\subsection{موارد موفق تشکیل کشور جدید}
\label{subsec:ch03-success}

\begin{table}[htbp]
	\centering
	\caption{موارد تشکیل کشور جدید و نتایج آن}
	\label{tab:ch03-new-states}
	\begin{adjustbox}{max width=\textwidth}
		\begin{tabular}{>{\bfseries\color{PrimaryMid}\small}p{2cm}
				>{\small}p{1.5cm} >{\small}p{2.5cm} >{\small}p{1.8cm}
				>{\small}p{4cm}}
			\toprule
			\textbf{کشور} & \textbf{سال} & \textbf{مبنا} &
			\textbf{پایداری} & \textbf{توضیح} \\
			\midrule
			بنگلادش &
			۱۹۷۱ &
			جنگ + مداخله‌ی هند &
			پایدار &
			ستم سیستماتیک پاکستان غربی \\
			\midrule
			اریتره &
			۱۹۹۳ &
			سی سال جنگ + همه‌پرسی &
			پایدار (ولی استبدادی) &
			استعمار اتیوپی \\
			\midrule
			تیمور شرقی &
			۲۰۰۲ &
			مداخله‌ی س.م. + همه‌پرسی &
			نسبتاً پایدار &
			اشغال اندونزی \\
			\midrule
			کوزوو &
			۲۰۰۸ &
			مداخله‌ی ناتو + اعلام یکجانبه &
			نسبتاً پایدار &
			شناسایی ناقص \\
			\midrule
			سودان جنوبی &
			۲۰۱۱ &
			توافق صلح + همه‌پرسی &
			ناپایدار &
			شکست دولت‌سازی \\
			\bottomrule
		\end{tabular}
	\end{adjustbox}
\end{table}

\subsection{موارد ناموفق یا ناتمام}
\label{subsec:ch03-failure}

\begin{tcolorbox}[enemybox,
	title={درس‌هایی از جدایی‌های ناتمام یا ناموفق}]
	\begin{itemize}[nosep, rightmargin=1em]
		\item \enemy{بیافرا (نیجریه، ۱۹۶۷--۷۰)}: جنگ داخلی با بیش
		از یک میلیون کشته؛ شکست جدایی؛
		\item \enemy{چچن (۱۹۹۴--۲۰۰۹)}: دو جنگ ویرانگر؛ شکست جدایی؛
		ویرانی گروزنی؛
		\item \enemy{تامیل ایلام (سریلانکا، ۱۹۸۳--۲۰۰۹)}: جنگ داخلی
		طولانی؛ شکست نظامی ببرهای تامیل؛
		\item \enemy{کردستان عراق (۲۰۱۷)}: همه‌پرسی با ۹۳ درصد آرای
		مثبت؛ مخالفت بغداد، ترکیه و ایران؛
		\item \enemy{کاتالونیا (۲۰۱۷)}: همه‌پرسی یکجانبه؛ سرکوب
		اسپانیا؛ بحران سیاسی بدون نتیجه.
	\end{itemize}
\end{tcolorbox}

\subsection{عوامل تعیین‌کننده‌ی موفقیت یا شکست}
\label{subsec:ch03-factors}

بر اساس مطالعه‌ی تطبیقی موارد فوق و پژوهش‌های
\thinker{رایان گریفیتس}%
\footnote{%
	رایان گریفیتس (\lr{Ryan D.\,Griffiths})؛%
	استاد علوم سیاسی دانشگاه سیراکیوز.%
}%
عوامل زیر در تعیین نتیجه‌ی جدایی نقش کلیدی
دارند\src{\lr{Griffiths, 2016, pp.\,15--43; Fazal, 2007, pp.\,1--28}}:

\begin{table}[htbp]
	\centering
	\caption{عوامل تعیین‌کننده‌ی موفقیت جدایی‌طلبی}
	\label{tab:ch03-factors}
	\begin{adjustbox}{max width=\textwidth}
		\begin{tabular}{>{\bfseries\color{PrimaryMid}}p{2.8cm}
				p{3.5cm} p{3.5cm} p{2cm}}
			\toprule
			\textbf{عامل} & \textbf{تسهیل‌کننده} &
			\textbf{بازدارنده} & \textbf{اهمیت} \\
			\midrule
			حمایت بین‌المللی &
			شناسایی قدرت‌های بزرگ &
			وتو در شورای امنیت &
			بسیار بالا \\
			\midrule
			ظرفیت دولت مرکزی &
			دولت ضعیف یا فروپاشیده &
			دولت قوی و متحد &
			بالا \\
			\midrule
			انسجام داخلی جنبش &
			رهبری یکپارچه &
			شکاف‌های درونی &
			بالا \\
			\midrule
			سرزمین مشخص &
			مرزهای تاریخی روشن &
			جمعیت درهم‌آمیخته &
			متوسط \\
			\midrule
			منابع اقتصادی &
			خودکفایی بالقوه &
			وابستگی به مرکز &
			متوسط \\
			\midrule
			هنجار بین‌المللی &
			سابقه‌ی جدایی مشابه &
			اصل تمامیت ارضی &
			بالا \\
			\bottomrule
		\end{tabular}
	\end{adjustbox}
\end{table}


% ────────────────────────────────────────────────────────────────
\section{نقد کلی: حقوق بین‌الملل، ابزار عدالت یا قدرت؟}
\label{sec:ch03-critique}

\begin{tcolorbox}[quotebox, title={}]
	«حقوق بین‌الملل، حقوقی است که اخلاق‌گرایان نادیده
	می‌گیرند و قدرتمندان تفسیر می‌کنند.»
	
	\hfill ---منتسب به
	ای.\,اچ.\,کار\src{\lr{Carr, 1939/2001, p.\,170}}
\end{tcolorbox}

تحلیل حقوقی این فصل نشان داد که:
\begin{enumerate}[label=\textbf{\arabic*.}, nosep, rightmargin=1em]
	\item حقوق بین‌الملل ابزارهایی هنجاری برای حمایت از مردمان
	تحت ستم فراهم کرده (حق تعیین سرنوشت، \lr{R2P}، کنوانسیون
	نسل‌کشی)؛
	\item اما اعمال این ابزارها کاملاً وابسته به \concept{اراده‌ی
		سیاسی} قدرت‌های بزرگ و ساختار شورای امنیت است؛
	\item جدایی جبرانی یک نظریه‌ی در حال تکوین است ولی هنوز به
	هنجار قطعی تبدیل نشده؛
	\item مرز حقوقی میان «حمایت از غیرنظامیان» و «تغییر رژیم»
	در عمل بسیار مبهم است؛
	\item حقوق بین‌الملل نه قادر به حل این تنش‌هاست و نه قادر
	به جلوگیری از سوءاستفاده؛ اما غیابش وضع را بدتر می‌کند.
\end{enumerate}

در فصل بعد، از حوزه‌ی حقوقی به \concept{حوزه‌ی تحلیلی-ارزیابی}
گذار می‌کنیم و الگوی چندبعدی ارزیابی مداخلات را تفصیل خواهیم داد.